% Part: first-order-logic
% Chapter: introduction
% Section: soundness-completeness

\documentclass[../../../include/open-logic-section]{subfiles}

\begin{document}

\olfileid{fol}{int}{scp}

\olsection{Kasahihan lan Kelengkapan}

Kita uga bakal ngenalake sistem !!{derivation} kanggo logika sepisanan. Ana
akeh sistem !!{derivation} kang wis dikembangake para ahli logika, nanging
kabeh nemtokake relasi !!{derivability} kang padha ing antarane
!!{sentence}s. Kita kandha yen $\Gamma$ \emph{!!{derive}s}~$!A$,
$\Gamma \Proves !A$, yen ana !!a{derivation} saka jinis tartamtu kang
ditemtokake kanthi persis. !!^{derivation}s mesthi arupa tatanan simbol
kang winates---mbokmenawa dhaftar !!{sentence}s, utawa sawijining
struktur kang luwih rumit. Tujuane sistem !!{derivation} yaiku nyawisake
piranti kanggo nemtokake apa !!a{sentence} dadi akibat saka sawijining
himpunan~$\Gamma$. Supaya bisa nggayuh tujuan iki, kudu lumaku yen
$\Gamma \Entails !A$ yen lan mung yen $\Gamma \Proves !A$.

Yen $\Gamma \Proves !A$ nanging ora $\Gamma \Entails !A$, sistem
!!{derivation} kita bakal kuwat banget lan mbuktekake kakehan. Sipat yen
$\Gamma \Proves !A$ menehi $\Gamma \Entails !A$ diarani
\emph{kasahihan}, lan iki syarat paling sithik kanggo saben sistem
!!{derivation} kang becik. Kosok baline, yen $\Gamma \Entails !A$
nanging ora $\Gamma \Proves !A$, sistem !!{derivation} kita ringkih
banget lan ora mbuktekake cukup akeh. Sipat yen $\Gamma \Entails !A$
menehi $\Gamma \Proves !A$ diarani \emph{kelengkapan}. Kasahihan
lumrahe rada gampang dibuktekake (kanthi induksi tumrap struktur
!!{derivation}s, kang ditemtokake kanthi induktif). Kelengkapan luwih
angel dibuktekake.

Kasahihan lan kelengkapan nduweni sawetara akibat wigati. Yen sawijining
himpunan !!{sentence}s~$\Gamma$ !!{derive}s kontradiksi (kayata
$!A \land \lnot !A$), himpunan kasebut diarani \emph{ora konsisten}.
$\Gamma$ kang ora konsisten ora bisa nduweni model apa wae; himpunan
kasebut ora bisa dicukupi. Saka kelengkapan, kosok baline ndherek: saben
$\Gamma$ kang ora kalebu ora konsisten---utawa, kaya kang bakal kita
kandhakake, \emph{konsisten}---nduweni model. Malah, iki padha karo
kelengkapan, lan wujud kelengkapan iki kang bakal kita buktekake. Iki
asil kang jero lan mbokmenawa nggumunake: kasunyatan yen sampeyan ora
bisa mbuktekake $!A \land \lnot !A$ saka $\Gamma$ wis njamin anane
!!a{structure} kang pas kaya katrangane~$\Gamma$. Dadi, kelengkapan
menehi wangsulan marang pitakon: himpunan !!{sentence}s endi kang nduweni
model? Wangsulane: persis kabeh himpunan kang konsisten.

Teorema kasahihan lan kelengkapan nduweni loro akibat wigati: teorema
kekompakan lan teorema L\"owenheim--Skolem. Iki asil-asiling wigati ing
teori model, lan bisa digunakake kanggo netepake akeh asil kang narik
kawigaten. Kita wis nyebut loro ing antarane: logika sepisanan ora bisa
ngandharake yen !!{domain} saka !!a{structure} iku winates utawa yen
struktur kasebut !!{nonenumerable}.

Ing sejarah, kabeh iki---carane nemtokake sintaksis lan semantik logika
sepisanan, carane nemtokake sistem !!{derivation} kang becik, carane
mbuktekake kasahihan lan kelengkapane, lan carane mangerteni kanthi
cetha apa kang bisa lan ora bisa diandharake ing basa logika
sepisanan---mbutuhake wektu suwe nganti ditemokake lan dirumusake kanthi
bener. Saiki kita wis ngerti carane nindakake, nanging ngliwati kabeh
rinciane isih bisa mbingungake lan marahi bosen. Nanging rincian kasebut
uga wigati, amarga metode kang dikembangake ing kene kanggo basa formal
logika sepisanan ditrapake ing endi-endi sajrone logika, ilmu komputer,
lan linguistik. Mula, nggarap rincian kasebut bakal migunani kanggo
jangka dawa.

\end{document}
